\documentclass[11pt]{article}
\usepackage[margin=1in]{geometry}
\usepackage{graphicx}
\usepackage{booktabs}
\begin{document}
% =========================
% TITLE PAGE
% =========================
\begin{titlepage}
\centering
\vspace*{2cm}
{\Huge \textbf{International Institute Of Information Technology Bangalore}}\\[1 cm]
{\Large AIM-102}\\[2cm]
\textbf{Name:} Om Tripathi\\[0.5cm]
\textbf{Roll Number:} BA2025027\\[0.5cm]
\textbf{Batch:} BT DSAI\\[1.5cm]
\textbf{Statistical Machine Learning}\\
\textbf{Course Project}\\
\vfill
\end{titlepage}
\newpage
% =========================
% INTRODUCTION
% =========================
\section{Introduction}
Natural Language Processing (NLP) is widely used for analyzing textual data. In this project, we compare classical machine learning models using two different feature representations: TF-IDF (sparse features) and contextual embeddings generated using DistilBERT.
The objective is to analyze how feature representation affects classification performance in sentiment analysis.
% =========================
% DATASET
% =========================
\section{Dataset Description and EDA}
We use the IMDb dataset from Hugging Face, which is a binary classification dataset for sentiment analysis.
\begin{itemize}
    \item Number of samples used: 2000 (subset)
    \item Classes: Positive and Negative
\end{itemize}
Basic exploratory data analysis (EDA) was performed, including class distribution and text length analysis.
% =========================
% METHODOLOGY
% =========================
\section{Methodology}
\subsection{TF-IDF Features}
Text preprocessing included cleaning and tokenization. TF-IDF vectorization was used to convert text into numerical feature vectors.
\subsection{Embedding Features}
We used DistilBERT as a pretrained encoder model to extract embeddings. The model was used strictly as a feature extractor without fine-tuning.
Mean pooling was applied to obtain sentence-level embeddings.
\subsection{Models}
The following classical models were used:
\end{document}